% Compile: pdflatex mcp-client.tex && pdflatex mcp-client.tex

\documentclass[11pt]{article}

\usepackage[margin=1in]{geometry}
\usepackage[utf8]{inputenc}
\usepackage[T1]{fontenc}
\usepackage{hyperref}
\usepackage{url}
\usepackage{booktabs}
\usepackage{amsmath}
\usepackage{amsfonts}
\usepackage{microtype}
\usepackage{enumitem}
\usepackage{graphicx}
\usepackage{xcolor}

\hypersetup{
  colorlinks=true,
  linkcolor=blue!60!black,
  citecolor=blue!60!black,
  urlcolor=blue!60!black,
}

\title{An MCP Client for XSEC: Attaching External Tool Servers\\
Without Bloating the Tool Surface or the Trust Boundary}

\author{
  XSEC Labs \\
  \texttt{github.com/uncesaii/xsec}
}

\date{}

\begin{document}

\maketitle

\begin{abstract}
The Model Context Protocol (MCP) lets an agent attach external tool servers ---
a filesystem server, a Burp or Nuclei bridge, a ticketing system, a cloud API
--- over a small JSON-RPC protocol. XSEC has no MCP client today, and a naive one
would fight two properties XSEC has deliberately built. First, MCP tool schemas
are re-sent every turn and a handful of servers can spend tens of thousands of
tokens on definitions before any work happens, undoing the 10--20-tool
always-loaded budget XSEC has just adopted. Second, and more seriously, an
attached server is an untrusted party: its results are attacker-influenceable
text fed straight to the model (prompt injection), and a remote server drives the
client to fetch URLs it chooses (SSRF, confused deputy, token passthrough). This
paper surveys how MCP works and how OpenCode and Claude Code expose it, then
designs an MCP-client layer built entirely from machinery XSEC already owns: an
\texttt{McpHost} modeled on the existing out-of-process \texttt{PluginHost};
server declarations in the two-level settings under a gated \texttt{Security}
group; discovered tools capability-gated through the same
\texttt{gateFlagsFor} path as plugins and \texttt{self\_extend}; remote egress
routed through the existing \texttt{ScopePolicy} / \texttt{PathPolicy} /
\texttt{RateLimiter} chokepoint; results neutralized by the already-MCP-aware
\texttt{untrusted-sanitizer}; and every MCP tool placed behind the deferred
loader, opt-in and off by default. The guiding principle: an MCP tool is an
externally-hosted tool, and XSEC already has a security model for that. This is a
design and decision log dated 2026-08-30. No code is written here.
\end{abstract}

\section{Motivation}

An autonomous pentesting agent benefits from tools its authors did not write:
an operator's private recon service, a licensed scanner behind an HTTP bridge,
a knowledge base of prior engagements, a bug-tracker to file findings into. MCP
is the emerging lingua franca for exactly this --- a standard client-server
protocol that a growing ecosystem of servers already speaks
\cite{mcp-arch}. Editors and agents (VS Code, Claude Code, OpenCode, Claude
Desktop) act as MCP \emph{hosts}; attaching a server is a config edit
\cite{mcp-arch,claudecode-mcp,opencode-mcp}.

XSEC is not an MCP host. It ships a fixed registry of model-facing tools
(\texttt{TOOL\_REGISTRY\_ORDER} in
\texttt{packages/core/src/agent/tools/index.ts}) assembled from per-domain
modules, and the only way to add a tool at runtime is \texttt{self\_extend},
which registers \emph{model-authored} tools into a single session
(\texttt{packages/core/src/agent/tools/system.ts}). There is no way to attach an
\emph{externally hosted} tool server. This paper asks what it would take to add
one without regressing the two things XSEC has spent effort getting right: a
small, well-selected tool surface, and a hard trust boundary around live
targets.

Those two things are in direct tension with a naive MCP integration, so they
frame the whole design:

\begin{itemize}[leftmargin=*]
  \item \textbf{Context cost.} MCP tool definitions are standing context. A
  setup wiring in a few servers can spend on the order of tens of thousands of
  tokens on tool schemas re-sent every turn, degrading both selection accuracy
  and cost \cite{mcp-clientbp,anthropic-toolsearch}. XSEC's companion
  tool-curation work \cite{tool-curation} sets a 10--20-tool budget per
  reasoning context and commits to a deferred loader; MCP must ride that, not
  bypass it.
  \item \textbf{Trust boundary.} An MCP server is an untrusted principal. Its
  tool results are attacker-influenceable text that flows straight into the
  model's context (prompt injection), and a remote server chooses the URLs the
  client fetches during discovery and auth (SSRF, confused deputy)
  \cite{mcp-security}. XSEC already treats live-target egress as a guarded
  chokepoint; MCP must be forced through the same gate, not around it.
\end{itemize}

\section{How MCP works, in brief}

MCP is a client-server protocol layered as a \emph{data layer} (a JSON-RPC 2.0
message protocol) over a \emph{transport layer} \cite{mcp-arch}. A host creates
one client per server and holds a dedicated connection to each.

\paragraph{Transports.} Two are current \cite{mcp-arch}:
\begin{itemize}[leftmargin=*]
  \item \textbf{stdio} --- the host spawns the server as a local child process
  and speaks over its standard input/output. Optimal for local servers, no
  network hop; but the server runs with the host's privileges (Section 7).
  \item \textbf{Streamable HTTP} --- HTTP POST for client-to-server messages
  with optional Server-Sent Events (SSE) for streaming. This is the remote
  transport and carries standard HTTP auth (bearer tokens, API keys, OAuth).
  The older pure-SSE transport is deprecated in favor of it
  \cite{claudecode-mcp}.
\end{itemize}

\paragraph{Discovery and invocation.} Over any transport the exchange is the
same JSON-RPC \cite{mcp-arch}:
\begin{enumerate}[leftmargin=*]
  \item An optional \texttt{server/discover} returns the server's identity,
  supported protocol versions, and capabilities (which primitives it offers,
  whether it emits change notifications).
  \item \texttt{tools/list} returns an array of tool descriptors, each with a
  \texttt{name}, \texttt{title}, \texttt{description}, and a JSON-Schema
  \texttt{inputSchema}. The host merges these into the tool set it advertises to
  the model.
  \item \texttt{tools/call} invokes a named tool with an \texttt{arguments}
  object and returns a \texttt{content} array (text, images, resource
  references). Tool errors arrive as a successful response with
  \texttt{isError: true} rather than as a transport failure
  \cite{mcp-clientbp}.
  \item Servers may push \texttt{notifications/tools/list\_changed} when their
  tool set changes; clients re-list in response \cite{mcp-arch}.
\end{enumerate}

Servers can also expose \emph{resources} (readable context) and \emph{prompts}
(templates); for a first XSEC integration only \emph{tools} matter, so the rest
of this paper scopes to tools.

\paragraph{Auth.} Remote servers authenticate with bearer tokens or API keys in
HTTP headers, and MCP recommends OAuth to obtain those tokens
\cite{mcp-arch,mcp-security}. The authorization spec forbids
\emph{token passthrough} --- a server accepting a token that was not issued for
it --- precisely because it breaks audience validation and enables confused-deputy
chains \cite{mcp-security}.

\section{Prior art: OpenCode and Claude Code}

Both of the agents XSEC most resembles already ship MCP clients, and their
choices are worth borrowing.

\paragraph{OpenCode.} Servers are declared under an \texttt{mcp} key in
\texttt{opencode.json}, keyed by name, each tagged \texttt{"local"} (with a
\texttt{command} array plus \texttt{cwd}/\texttt{environment}/\texttt{timeout}) or
\texttt{"remote"} (with a \texttt{url} plus \texttt{headers}/\texttt{oauth}), and
each with an explicit \texttt{enabled} flag \cite{opencode-mcp}. Tools can be
gated with glob patterns in a \texttt{tools} map (e.g. \texttt{"my-mcp*": false}),
globally or per-agent, and remote OAuth is driven by
\texttt{opencode mcp auth <server>}. The docs carry the blunt warning that
motivates half of this paper: \emph{``MCP servers add to your context, so you
want to be careful with which ones you enable''} \cite{opencode-mcp}.

\paragraph{Claude Code.} Servers live in \texttt{.mcp.json} (project scope,
checked into git) or in the user's global config, keyed under
\texttt{mcpServers} \cite{claudecode-mcp}. Transports are \texttt{stdio},
\texttt{http} (recommended), deprecated \texttt{sse}, and \texttt{ws}. Servers
are added with \texttt{claude mcp add}, e.g.

\begin{quote}
\ttfamily
claude mcp add --transport http github --scope project https://.../mcp/\\
claude mcp add --transport stdio airtable -- npx -y airtable-mcp-server
\end{quote}

Three lessons carry directly into the XSEC design:
\begin{itemize}[leftmargin=*]
  \item \textbf{Namespacing.} MCP tools are surfaced to the model as
  \texttt{mcp\_\_<server>\_\_<tool>}, so a server's tools cannot collide with
  built-ins and are addressable in permission rules by that full name
  \cite{claudecode-mcp}.
  \item \textbf{Consent and trust.} Project-scoped servers require explicit
  approval on first use; trust is cached per folder. The docs warn plainly that
  \emph{servers that fetch external content can expose you to prompt injection}
  and that you should verify you trust each server before connecting
  \cite{claudecode-mcp}.
  \item \textbf{Output limits.} MCP tool output is capped (a 10K-token warning
  threshold, a 25K default limit, tunable via \texttt{MAX\_MCP\_OUTPUT\_TOKENS}),
  acknowledging that tool \emph{results}, not just schemas, are a context tax
  \cite{claudecode-mcp}.
\end{itemize}

\section{The context-cost problem and deferred loading}

The MCP client best-practices document is explicit that loading every tool
upfront ``breaks down'' as servers accumulate: definitions alone can consume the
majority of the context window before the model reads the user's message
\cite{mcp-clientbp}. Its recommended fix is \emph{progressive discovery}: fetch
definitions via \texttt{tools/list} but defer injecting them, expose a
lightweight \texttt{search\_tools} meta-tool, and load full schemas only when the
model asks. It recommends switching to this mode once tool definitions exceed
roughly 1--5\% of the context window \cite{mcp-clientbp}.

This is the same conclusion XSEC's tool-curation work reached from the other
direction \cite{tool-curation}: reliable selection lives in a 10--20 always-loaded
tool band, accuracy collapses non-linearly past 30--50 tools, and the fix is a
deferred \texttt{list\_tools} / \texttt{load\_tool} layer mirroring the existing
JIT-skills pattern (\texttt{list\_skills} / \texttt{load\_skill}, gated by
\texttt{jitSkills}, in \texttt{packages/core/src/agent/tools/skills.ts}). The two
documents agree on the mechanism and the numbers.

The implication for MCP is unambiguous: \textbf{MCP tools must not be
always-loaded.} A server may expose dozens of tools; even one GitHub-class
server can dominate the surface. MCP tools belong in the long tail, behind the
deferred loader, so the base context the model sees on every turn stays inside
the band. The MCP client caches each server's \texttt{tools/list} result
host-side (respecting the \texttt{ttlMs} / \texttt{list\_changed} freshness
hints \cite{mcp-clientbp}) and indexes it for the loader; nothing enters model
context until the model searches for and loads a specific tool. This also plays
well with prompt caching --- appending a newly loaded definition after the cache
breakpoint, rather than re-sorting the whole \texttt{tools} array, avoids
invalidating the cached prefix \cite{mcp-clientbp}.

\section{XSEC today: the parts we build on}

The integration is far cheaper than it looks because XSEC already contains a
near-exact blueprint: an out-of-process, capability-gated tool \emph{provider}.
Six pieces do the work.

\paragraph{The plugin system --- the blueprint.} \texttt{packages/core/src/plugins/}
already implements a \texttt{PluginHost} that runs an external subprocess
speaking JSON over stdio, discovers its tools (\texttt{list\_tools}), invokes
them (\texttt{call\_tool}), and contributes their definitions into the model's
tool set through the single authorization path (protocol in
\texttt{plugins/protocol.ts}; design in \texttt{plugins/DESIGN.md}). An MCP client
is architecturally the \emph{same shape} --- a subprocess or remote endpoint that
contributes gated tool definitions --- so it should be built as an
\texttt{McpHost} sibling of \texttt{PluginHost}, not as a new subsystem. The
\texttt{PluginHost} interface (\texttt{plugins/loader.ts}) is the contract to
mirror: \texttt{toolDefinitions()}, \texttt{gateMaps()},
\texttt{capabilityFlagsFor(name)}, \texttt{ownsTool(name)},
\texttt{call(name, args)}. Notably, the \texttt{@modelcontextprotocol/sdk}
dependency is \emph{already present}: \texttt{packages/cli/src/commands/mcp-server.ts}
uses it to expose XSEC's own tools \emph{as} an MCP server (behind an
\texttt{MCP\_LIVE\_TOOL\_NAMES} allowlist). The client proposed here is the
inverse direction and can share the SDK and naming conventions already
established there.

\paragraph{The tool registry.} \texttt{TOOL\_REGISTRY\_ORDER}
(\texttt{tools/index.ts}) enumerates the built-in tools (57 entries; the console
role advertises a $\sim$38-tool default subset); per-domain modules each own
their definitions and a co-located dispatch map. \texttt{ToolDefinition}
(\texttt{agent/types.ts}) is converted to the Anthropic native schema by
\texttt{toNativeToolDef} / \texttt{toNativeExtensionToolDef}
(\texttt{agent/native-tooldef.ts}). Crucially, the interactive console already
has an injection seam: \texttt{refreshInjectedTools()} in
\texttt{packages/core/src/console/turn-engine.ts} runs at each turn boundary and
merges \texttt{config.pluginHost.toolDefinitions()} and its gate maps into the
advertised set. An \texttt{McpHost} folds into that same seam via a new
\texttt{config.mcpHost?}, which is what makes it reuse the authorization path for
free.

\paragraph{The capability model.} XSEC already admits externally-authored tools
at runtime, and the capability vocabulary is a closed set in
\texttt{plugins/manifest.ts}: \texttt{PluginCapability} is one of
\texttt{network}, \texttt{filesystem-read}, \texttt{filesystem-write},
\texttt{process-exec}, \texttt{findings-write}. \texttt{gateFlagsFor(tool)} is the
single conservative translation from declared capabilities to the runtime gates
(\texttt{network} or \texttt{process-exec} implies network-capable;
\texttt{filesystem-*} implies local-scope; read-only only when every capability
is a pure read; unknown maps to the most restrictive), and
\texttt{validatePluginManifest} makes capabilities mandatory and non-empty,
rejects built-in name collisions, and fails closed. The gates themselves are
keyed on tool \emph{name} through three maps in
\texttt{console/turn-engine.ts} --- \texttt{NETWORK\_CAPABLE\_TOOLS},
\texttt{LOCAL\_SCOPE\_TOOLS}, \texttt{READ\_ONLY\_TOOLS} --- and
\texttt{plugins/DESIGN.md} states the hazard bluntly: a tool absent from all
three is treated as the least-dangerous class (``danger by omission of
safety''). This is the exact hole an MCP client must not open. The
\texttt{self\_extend} tool (\texttt{tools/system.ts}, handler
\texttt{ToolExecutor.selfExtend}, registry \texttt{plugins/self-extension.ts}) is
the same model for \emph{model-authored} tools --- additive-only, session-scoped,
off unless \texttt{allowModelSelfExtension} is set, capabilities mandatory, with
per-session limits (8 extensions, 32 tools, 16KiB/manifest) --- but with one
telling difference: for \texttt{self\_extend}, registration is \emph{not}
execution (a self-authored tool has no in-process body). MCP is the case where we
\emph{do} execute the external tool, which is exactly why routing it through the
same capability gate matters more, not less.

\paragraph{The egress / scope chokepoint.} Live-target traffic is guarded by a
stack under \texttt{packages/core/src/scope/}: a host allowlist
(\texttt{ScopePolicy}), a path-prefix allowlist (\texttt{PathPolicy}), a per-host
\texttt{RateLimiter}, and an \texttt{EnforcementTracker} threaded from
\texttt{AgentConfig} down through \texttt{ToolContext} so ``every fetch
chokepoint mutates the same counters''
(\texttt{scope/enforcement.ts}). \texttt{scope-guard.ts} makes the absence of a
scope \emph{visible} (emitting a \texttt{scope\_guards\_inert} event) and
\texttt{0SEC\_REQUIRE\_SCOPE} makes every mode fail closed;
\texttt{agenticScan()} refuses a live network target before initializing any
client when no scope is configured. This is the machinery a remote MCP server's
egress must be forced through.

\paragraph{The untrusted-input sanitizer --- already MCP-aware.} The single most
convenient fact for this design: \texttt{packages/core/src/untrusted-sanitizer.ts}
already anticipates MCP. \texttt{isUntrustedSourceTool(name)} returns true for
network/read tools \emph{and for any name beginning with} \texttt{mcp\_\_}, and
\texttt{sanitizeUntrustedToolResult} deterministically neutralizes injection
patterns, wrapping output in \texttt{[[0SEC\_UNTRUSTED\_DATA]]} with a
DATA-not-instructions note. It is already wired into
\texttt{agent/native-loop.ts} (which emits an \texttt{untrusted\_input\_sanitized}
event), its module header explicitly lists ``MCP tool outputs'' as a target, and
\texttt{PluginHost} already routes results through it. An MCP client inherits this
defense for free simply by using the \texttt{mcp\_\_<server>\_\_<tool>} naming
convention.

\paragraph{The config surface.} Configuration is a two-level settings file: a
per-user global (\texttt{\textasciitilde/.xsec/tui-settings.json}) with a
per-project override (\texttt{<cwd>/.xsec/tui-settings.json}), described by
\texttt{SETTING\_DEFS} in \texttt{packages/cli/src/tui/settings.ts} and operated
by \texttt{xsec config show/export/import}
(\texttt{packages/cli/src/commands/config.ts}). Settings carry a \texttt{group}
label, and keys in the \texttt{Security} group (e.g.
\texttt{allowModelSelfExtension}) get special handling: \texttt{config import}
refuses to silently enable a Security-group key unless \texttt{--yes} is passed,
prints the specific change, drops unknown keys, and validates every value
through a sanitiser before it lands on disk. Feature flags themselves are
env-driven getters in \texttt{packages/core/src/agent/features.ts}
(\texttt{0SEC\_FEATURE\_*}), most defaulting off.

\section{Design}

The design principle in one line: \textbf{an MCP tool is a model-facing,
externally-hosted tool, and XSEC already has a security model for
externally-authored tools --- reuse it, do not invent a second one.}

\subsection{Where servers are declared}

Add an \texttt{mcp} block to the two-level settings, following OpenCode's shape
but living inside XSEC's existing layered config so global/project precedence and
the import-consent gate apply for free:

\begin{quote}
\ttfamily
"mcp": \{\\
\ \ "recon-kb": \{ "transport": "stdio", "command": ["npx","-y","kb-mcp"],\\
\ \ \ \ "env": \{"KB\_TOKEN":"\$\{KB\_TOKEN\}"\}, "enabled": false,\\
\ \ \ \ "capabilities": ["network","findings-write"],\\
\ \ \ \ "tools": \{ "search*": true, "delete*": false \} \},\\
\ \ "burp": \{ "transport": "http", "url": "https://127.0.0.1:1337/mcp",\\
\ \ \ \ "headers": \{"Authorization":"Bearer \$\{BURP\_KEY\}"\},\\
\ \ \ \ "enabled": false, "capabilities": ["network"],\\
\ \ \ \ "scope\_bound": true \} \}
\end{quote}

Per-server fields: \texttt{transport} (\texttt{stdio} or \texttt{http}),
\texttt{command}/\texttt{env} or \texttt{url}/\texttt{headers}, an explicit
\texttt{enabled} defaulting to \texttt{false}, a per-server \texttt{capabilities}
ceiling (see 6.3), a \texttt{tools} glob allow/deny map (borrowed from OpenCode),
and \texttt{scope\_bound} declaring that a remote server's egress must satisfy
the active \texttt{ScopePolicy}. The whole \texttt{mcp} block is a
\texttt{Security}-group key, so importing a shared config can never silently
attach a server; and a new master flag \texttt{0SEC\_FEATURE\_MCP\_CLIENT}
(default off) gates the subsystem entirely, matching how every other sharp
feature ships.

\subsection{Transports}

Support \texttt{stdio} and streamable \texttt{http} first; pure SSE is deprecated
upstream \cite{claudecode-mcp} and can wait. stdio servers are local child
processes --- treat their \texttt{command} exactly as OpenCode/Claude Code do,
but spawn them with the same hardened child environment XSEC already builds for
plugins (\texttt{buildPluginEnv} over \texttt{allowlistedChildEnv} in
\texttt{agent/sanitized-env.ts}, which never inherits provider/cloud credentials
or the target-auth block), and see Section 7 for the consent requirement before
spawning. http servers get
their egress routed through the scope chokepoint (6.4). Change notifications
(\texttt{tools/list\_changed}) are honored by re-listing and re-indexing, but
treated as a catalog-refresh event, not a per-turn one, to protect prompt
caching \cite{mcp-clientbp}.

\subsection{How discovered tools enter the registry}

On connect, the client runs \texttt{server/discover} then \texttt{tools/list},
applies the server's \texttt{tools} glob filter, and produces registry entries
named \texttt{mcp\_\_<server>\_\_<tool>} --- the Claude Code convention
\cite{claudecode-mcp} --- so MCP tools never collide with XSEC built-ins and are
addressable by full name in any allow/deny rule. Each entry's
\texttt{inputSchema} becomes its parameter schema; its dispatch route points at a
single generic executor that marshals \texttt{tools/call} to the owning client.

Three constraints on admission:
\begin{enumerate}[leftmargin=*]
  \item \textbf{Deferred, never always-loaded.} MCP tool descriptors are
  registered and indexed but \emph{not} injected into any per-role base set. They
  are reachable only through the deferred \texttt{list\_tools} /
  \texttt{load\_tool} front door \cite{tool-curation}, so the always-loaded band
  stays at 10--20 built-ins regardless of how many servers are attached. This is
  the single most important design choice for context hygiene.
  \item \textbf{Capability-bounded.} Every MCP tool inherits the owning server's
  declared \texttt{capabilities} ceiling, run through \texttt{gateFlagsFor} so it
  lands in the same \texttt{NETWORK\_CAPABLE\_TOOLS} / \texttt{LOCAL\_SCOPE\_TOOLS}
  / \texttt{READ\_ONLY\_TOOLS} maps as every other gated tool. MCP schemas carry
  no XSEC capabilities, so the client must \emph{assign} them, and to avoid the
  ``danger by omission'' hole every MCP tool defaults to at least
  \texttt{network} (it egresses to a server by definition) --- never the
  least-dangerous class by accident. A server declared \texttt{["network"]}
  cannot yield a tool that writes the filesystem, regardless of its schema.
  \item \textbf{Output-bounded.} Tool results are truncated to a configurable
  token budget (Claude Code's \texttt{MAX\_MCP\_OUTPUT\_TOKENS} pattern
  \cite{claudecode-mcp}) before they reach the model, bounding the result-side
  context tax and the injection payload size.
\end{enumerate}

\subsection{Fitting the security posture}

Three XSEC chokepoints do the heavy lifting:

\begin{itemize}[leftmargin=*]
  \item \textbf{Capability gate = the \texttt{self\_extend} gate.} MCP tools
  route through the same authorization check that governs model-authored tools:
  the declared capability set is the ceiling, enforced at dispatch, never
  self-expandable. No new authorization concept is introduced.
  \item \textbf{Remote egress = the scope chokepoint.} An \texttt{http} server's
  connection, discovery fetches, and every \texttt{tools/call} round trip are
  issued through the same fetch path that \texttt{ScopePolicy} /
  \texttt{PathPolicy} / \texttt{RateLimiter} / \texttt{EnforcementTracker} guard,
  so an out-of-scope MCP endpoint is refused and counted like any other egress.
  When \texttt{scope\_bound} is set and no scope is configured, the server fails
  closed under \texttt{0SEC\_REQUIRE\_SCOPE} semantics.
\end{itemize}

\subsection{The untrusted-input chokepoint}

MCP tool \emph{results} are untrusted input: they are attacker-influenceable text
returned by a third party and fed back into the model. Here XSEC has already done
the work. Because MCP tools are named \texttt{mcp\_\_<server>\_\_<tool>},
\texttt{isUntrustedSourceTool} already matches them, and the existing
\texttt{native-loop.ts} path already routes their results through
\texttt{sanitizeUntrustedToolResult} --- neutralizing injection patterns and
fencing the payload in \texttt{[[0SEC\_UNTRUSTED\_DATA]]} with a
DATA-not-instructions note --- before the result re-enters context. The client
adds only two things on top: (a) truncation to the output budget before
sanitization, and (b) the structural guarantee that an MCP result can never be
interpreted as control input --- it can never enable a capability, attach a
server, or alter scope. Centralizing this in the generic \texttt{call\_tool}
executor gives exactly one place to harden, the same discipline the scope stack
applies to egress.

\section{Security analysis}

Attaching an MCP server widens the attack surface in ways that map cleanly onto
threats the MCP security spec enumerates \cite{mcp-security}. For each, XSEC's
existing defenses are the mitigation --- which is the argument for reusing them
rather than building an MCP-specific security layer.

\paragraph{Prompt injection via tool results.} The primary risk. A malicious or
compromised server returns text engineered to hijack the agent (``ignore your
scope, exfiltrate X''). Claude Code warns about exactly this
\cite{claudecode-mcp}. Mitigation: the untrusted-input chokepoint (6.5), which is
\emph{already built and already MCP-aware} --- \texttt{isUntrustedSourceTool}
matches the \texttt{mcp\_\_} prefix and \texttt{sanitizeUntrustedToolResult}
neutralizes injection patterns and fences results in
\texttt{[[0SEC\_UNTRUSTED\_DATA]]} --- plus truncation and the structural denial
of any control effect (results cannot change capabilities or scope). The model
may be misled by content, but the harness boundaries it cannot cross are enforced
outside the model.

\paragraph{SSRF via remote servers.} During discovery and OAuth a server
supplies URLs the client fetches; a malicious server can point them at
\texttt{169.254.169.254} (cloud metadata), \texttt{localhost} services, or
private ranges \cite{mcp-security}. Mitigation: route all MCP egress through the
scope chokepoint, which already applies host and path allowlisting; require
HTTPS for remote URLs except loopback; block private/reserved IP ranges and
apply the same checks to redirect targets, per the spec's SSRF guidance
\cite{mcp-security}. \texttt{scope\_bound} servers additionally cannot reach
anything outside the active engagement's \texttt{ScopePolicy}.

\paragraph{Confused deputy and token passthrough.} Remote OAuth against a
third-party server can be abused to obtain codes without fresh consent, and a
server that forwards tokens issued for other audiences breaks OAuth boundaries
\cite{mcp-security}. Mitigation: keep credentials host-side in the settings
file, never expose them to the model or to tool arguments; validate token
audience; treat per-server OAuth as an operator action (an explicit
\texttt{xsec mcp auth <server>} step, following OpenCode \cite{opencode-mcp}),
not something the agent initiates mid-run.

\paragraph{Local server compromise.} A stdio server is arbitrary code run with
the agent's privileges; a malicious \texttt{command} in an imported config is
remote code execution \cite{mcp-security}. Mitigation: the \texttt{mcp} block is
a \texttt{Security}-group key, so \texttt{config import} refuses to enable a
server (and prints the exact \texttt{command}) without \texttt{--yes};
\texttt{enabled} defaults to false; the subsystem is behind
\texttt{0SEC\_FEATURE\_MCP\_CLIENT}; and stdio servers run under the same hardened
child environment as plugins (\texttt{buildPluginEnv} / \texttt{allowlistedChildEnv},
no inherited secrets or target auth), with capability declarations bounding what
the resulting tools may do.

\paragraph{Surface and scope-minimization.} Over-broad capability grants inflate
blast radius \cite{mcp-security}. Mitigation: the per-server \texttt{capabilities}
ceiling and per-tool glob gating enforce least privilege; deferred loading keeps
un-needed tools out of context entirely; and no MCP capability is granted that
the operator did not declare.

\paragraph{Residual risk.} A server the operator explicitly trusts, enables, and
grants \texttt{network} can still return misleading content within its
capability, and a tool the model deliberately loads and calls will run. XSEC's
defenses bound \emph{authority}, not \emph{truthfulness}; the untrusted-content
framing and human-in-the-loop enable/consent steps are what keep a bad server
from silently escalating. This is the same residual-trust posture that governs
\texttt{self\_extend} and live-target scanning today.

\section{Decision}

We will:
\begin{enumerate}[leftmargin=*]
  \item \textbf{Add an opt-in MCP client} behind \texttt{0SEC\_FEATURE\_MCP\_CLIENT}
  (default off), built as an \texttt{McpHost} sibling of the existing
  \texttt{PluginHost} and folded into \texttt{refreshInjectedTools()} via a new
  \texttt{config.mcpHost?}, supporting \texttt{stdio} and streamable \texttt{http}
  transports and sharing the already-present \texttt{@modelcontextprotocol/sdk}.
  \item \textbf{Declare servers in the existing two-level settings} under an
  \texttt{mcp} block that is a \texttt{Security}-group key, so the
  global/project precedence and the import-consent gate apply, and every server
  defaults to \texttt{enabled: false}.
  \item \textbf{Admit discovered tools deferred, never always-loaded.} MCP tools
  are namespaced \texttt{mcp\_\_<server>\_\_<tool>}, indexed for the deferred
  \texttt{list\_tools} / \texttt{load\_tool} front door \cite{tool-curation}, and
  kept out of every per-role base set so the always-loaded band stays 10--20.
  \item \textbf{Reuse the \texttt{self\_extend} capability gate} for MCP tools:
  each server declares a \texttt{capabilities} ceiling, enforced at dispatch,
  never self-expandable.
  \item \textbf{Route all remote MCP egress through the scope chokepoint}
  (\texttt{ScopePolicy} / \texttt{PathPolicy} / \texttt{RateLimiter} /
  \texttt{EnforcementTracker}), with HTTPS-only and private-range blocking for
  discovery/auth URLs, and \texttt{scope\_bound} servers failing closed under
  \texttt{0SEC\_REQUIRE\_SCOPE}.
  \item \textbf{Fence MCP tool results as untrusted content} at one chokepoint in
  the generic \texttt{tools/call} executor: truncate to an output budget, label
  provenance, and deny any control effect.
\end{enumerate}

\textbf{Why.} MCP is worth having --- it lets operators extend the agent without
core-team shipping --- but a naive client would regress both the tool-surface
discipline we just established \cite{tool-curation} and the trust boundary the
scope stack enforces. Every requirement above is satisfied by machinery XSEC
already owns: an out-of-process capability-gated tool provider
(\texttt{PluginHost}) whose shape the client copies, a shared capability
vocabulary and gate translation (\texttt{plugins/manifest.ts},
\texttt{gateFlagsFor}), an already-MCP-aware untrusted-result sanitizer
(\texttt{untrusted-sanitizer.ts}), a guarded egress chokepoint (\texttt{scope/}),
a consent-gated config layer (\texttt{Security}-group import), and a deferred tool
loader (in progress) --- plus the \texttt{@modelcontextprotocol/sdk} already on
disk. The integration is mostly \emph{wiring MCP into existing gates}, which is
exactly why it is safe to do.

\section{Phased plan}

\begin{enumerate}[leftmargin=*]
  \item \textbf{Phase 0 --- read-only stdio spike.} Behind the master flag,
  connect to a single local stdio server, run \texttt{server/discover} and
  \texttt{tools/list}, and index the results. No invocation yet. Validates the
  client, the descriptor-to-registry mapping, and the namespacing.
  \item \textbf{Phase 1 --- deferred invocation.} Wire MCP tools into the
  \texttt{list\_tools} / \texttt{load\_tool} loader and the generic
  \texttt{tools/call} executor, with the untrusted-result chokepoint and output
  budget. Enforce the per-server capability ceiling at dispatch. Confirm the
  always-loaded band is unchanged with servers attached.
  \item \textbf{Phase 2 --- config and consent.} Land the \texttt{mcp}
  settings block as a \texttt{Security}-group key, the \texttt{enabled}-default-off
  behavior, the \texttt{tools} glob gating, and the \texttt{config import}
  consent path (print \texttt{command}/\texttt{url}, require \texttt{--yes}).
  \item \textbf{Phase 3 --- remote HTTP + scope binding.} Add the streamable
  HTTP transport with egress routed through the scope chokepoint, HTTPS-only and
  private-range blocking on discovery/auth URLs, and \texttt{scope\_bound}
  fail-closed semantics.
  \item \textbf{Phase 4 --- remote auth.} Add \texttt{xsec mcp auth <server>} for
  operator-driven OAuth/token setup, credentials stored host-side and never
  exposed to the model, with audience validation and no token passthrough.
  \item \textbf{Phase 5 --- lifecycle polish.} Honor \texttt{tools/list\_changed}
  as a catalog-refresh (cache-friendly), add \texttt{xsec mcp list/get/remove}
  management commands, and surface per-server enforcement counters in the run
  log alongside the existing scope counters.
\end{enumerate}

\section{Open questions and tradeoffs}

\begin{itemize}[leftmargin=*]
  \item \textbf{Capability inference.} A server declares one ceiling, but its
  tools may vary (a read tool and a write tool under one \texttt{network} grant).
  Do we require per-tool capability annotations, or accept the server-level
  ceiling as the coarse bound? The coarse bound is simpler and matches
  \texttt{self\_extend}'s per-manifest granularity; per-tool refinement can come
  later.
  \item \textbf{Deferred-loading round trip.} Placing MCP tools behind the loader
  costs a \texttt{list\_tools} $\rightarrow$ \texttt{load\_tool} $\rightarrow$
  use hop before a cold MCP tool can run \cite{tool-curation}.
  That is the right trade for the long tail; if a specific server proves
  hot-path for a workflow, a per-phase loadout can promote a named subset, but
  the default must never be always-on.
  \item \textbf{Resources and prompts.} This design scopes to MCP \emph{tools}.
  Resources (readable context) and prompts (templates) are real primitives
  \cite{mcp-arch} but each is its own context-cost and trust question; defer
  until tools are proven.
  \item \textbf{Result provenance in evidence.} If an MCP tool contributes to a
  finding, its untrusted provenance should be recorded in the evidence lineage,
  so a reviewer can tell that part of a finding rests on third-party output. This
  intersects the report writer and is out of scope here.
\end{itemize}

\begin{thebibliography}{99}

\bibitem{mcp-arch}
Model Context Protocol. \textit{Architecture overview}.
\url{https://modelcontextprotocol.io/docs/concepts/architecture}

\bibitem{mcp-clientbp}
Model Context Protocol. \textit{Client best practices: progressive tool
discovery and programmatic tool calling}.
\url{https://modelcontextprotocol.io/docs/develop/clients/client-best-practices}

\bibitem{mcp-security}
Model Context Protocol. \textit{Security best practices} (attack vectors:
confused deputy, token passthrough, SSRF, local server compromise, scope
minimization).
\url{https://modelcontextprotocol.io/specification/2026-07-28/basic/security_best_practices}

\bibitem{opencode-mcp}
OpenCode. \textit{MCP servers}.
\url{https://opencode.ai/docs/mcp-servers/}

\bibitem{claudecode-mcp}
Claude Code. \textit{Connect Claude Code to tools via MCP}.
\url{https://code.claude.com/docs/en/mcp}

\bibitem{anthropic-toolsearch}
Anthropic. \textit{Tool search tool}. Claude Platform Documentation.
\url{https://platform.claude.com/docs/en/agents-and-tools/tool-use/tool-search-tool}

\bibitem{tool-curation}
XSEC Labs. \textit{Tool Curation for Security Agents}.
\texttt{research/papers/tool-curation/tool-curation.tex}, 2026.

\end{thebibliography}

\end{document}
